\documentclass[11pt,a4paper,oneside]{book}
\usepackage{90-faelles/latex/studiearkiv}
\hypersetup{pdftitle={Matematik – Bind 6: Variationsregning og symmetrier}}
\begin{document}
\frontmatter
\begin{titlepage}
\centering
\vspace*{30mm}
{\sffamily\color{arkivblaa}\Large MATEMATIK TIL UNIVERSITETSFYSIK\par}
\vspace{16mm}
{\Huge\bfseries Bind 6\\[6mm]Variationsregning og symmetrier\par}
\vspace{13mm}
{\Large Formler, begreber og matematisk sammenhæng\par}
\vfill
\begin{minipage}{0.8\textwidth}
Variationsprincipper, Noethers sætning, gruppestrukturer og symmetrier.

Dette bind genbruger de samme LaTeX-kapitler som det samlede matematikkompendium. Planlagte kapitler er markeret som ufærdige. Der indgår ikke separate gennemregnede opgaveeksempler.
\end{minipage}
\vspace{14mm}
{\small Fysik-studiearkiv · Arbejdsudgave\par}
\end{titlepage}
\chapter*{Læsevejledning}
\addcontentsline{toc}{chapter}{Læsevejledning}
Dette bind er en emneopdelt udgave af det samlede matematikkompendium. Den generelle kilde og alle figurer vedligeholdes centralt under \texttt{00-matematik/chapters} og \texttt{00-matematik/figures}. For hele overblikket samt henvisninger mellem bindene anvendes den samlede udgave. Ufuldstændige kapitler er tydeligt markerede.
\begingroup\small\tableofcontents\endgroup
\mainmatter
\chapter{Variationsregning}
\planlagt{Definitioner, notation, forudsætninger, centrale resultater og fysikforbindelser udvikles samlet.}
\section{Funktionaler og første variation}
\leavevmode\par
% Redaktionelt: tilføj definitioner, gyldighed, relationer og relevante figurer.

\section{Euler–Lagrange og randbetingelser}
\leavevmode\par
% Redaktionelt: tilføj definitioner, gyldighed, relationer og relevante figurer.

\section{Bivariable, bindinger og multiplikatorer}
\leavevmode\par
% Redaktionelt: tilføj definitioner, gyldighed, relationer og relevante figurer.

\section{Symmetrier og Noethers sætning}
\leavevmode\par
% Redaktionelt: tilføj definitioner, gyldighed, relationer og relevante figurer.

\chapter{Gruppeteori og symmetrier}
\planlagt{Definitioner, notation, forudsætninger, centrale resultater og fysikforbindelser udvikles samlet.}
\section{Grupper, homomorfier og repræsentationer}
\leavevmode\par
% Redaktionelt: tilføj definitioner, gyldighed, relationer og relevante figurer.

\section{Liegrupper, Liealgebraer og generatorer}
\leavevmode\par
% Redaktionelt: tilføj definitioner, gyldighed, relationer og relevante figurer.

\section{Rotationer, SO(3) og SU(2)}
\leavevmode\par
% Redaktionelt: tilføj definitioner, gyldighed, relationer og relevante figurer.

\section{Symmetrier og bevarelseslove}
\leavevmode\par
% Redaktionelt: tilføj definitioner, gyldighed, relationer og relevante figurer.

\end{document}
